% GENERATED FILE. Edit canonical lessons, then run npm run generate:books.
% canonical-source-hash: fnv1a64:db783de4310fd915
% canonical-lessons: LA-C368-orator, LA-C368-eloquentia, LA-C368-disputatio, LA-C368-dissero, LA-C368-eloquor

\chapter{The Orator}
\label{ch:la-the-orator}
\hlchaptermodality{368}

\begin{chapteropening}
\textbf{By the end of this chapter:} \emph{I can talk about the orator and his eloquence, a debate, and holding forth on a subject.}

Complete the last lesson of chapter 368: I can talk about the orator and his eloquence, a debate, and holding forth on a subject.
\end{chapteropening}

\section[ōrātor, ōrātōris]{ōrātor, ōrātōris — a speaker, an orator; an envoy}

\label{lesson:LA-C368-orator}



\begin{quote}
Before the new one: say the Latin for that … not, then the Latin for as.
\end{quote}

\subsection*{You'll want to know: ōrātor, ōrātōris}
\textbf{ōrātor, ōrātōris} — \textquotedblleft{}a speaker, an orator; an envoy\textquotedblright{}.

\begin{quote}
Ōrātor ā populō laudātus est.
\end{quote}

\begin{quote}
\emph{The speaker was praised by the people.}
\end{quote}

\emph{ōrātor, ōrātōris} m. Built on \emph{ōrō}, to speak, to plead. An envoy who speaks for his city is also an \emph{ōrātor}. English \emph{orator}.

\begin{cousinweb}
\emph{ōrātor} is built on \emph{ōrō} \textquotedblleft{}to speak, to plead\textquotedblright{} (Lewis \& Short), with the ending \emph{-tor} of the one who does it. English \emph{orator} comes from it.
\end{cousinweb}

\begin{grammarlens}[title={What you've built}]
One more word for telling, explaining and describing at B1.
\end{grammarlens}

\subsection*{Your turn}
\begin{itemize}
\raggedright
  \item \emph{Say it:} \emph{ōrātor, ōrātōris}
  \item \emph{Say it:} \emph{ōrātor, ōrātōris}, once more
  \item \emph{Say it:} a sentence of your own with \emph{ōrātor, ōrātōris}
  \item \emph{Recall:} say the Latin for that … not, then the Latin for as, then say \emph{ōrātor, ōrātōris} again
\end{itemize}

\begin{tcolorbox}[breakable,colback=teal!4,colframe=teal!35!black,title={Before you move on}]
What does ōrātor, ōrātōris mean? (\textquotedblleft{}A speaker, an orator; an envoy\textquotedblright{}.) Say it once more.
\end{tcolorbox}
\section[ēloquentia, ēloquentiae]{ēloquentia, ēloquentiae — eloquence}

\label{lesson:LA-C368-eloquentia}



\begin{quote}
Before the new one: say the Latin for as, then the Latin for a speaker.
\end{quote}

\subsection*{You'll want to know: ēloquentia, ēloquentiae}
\textbf{ēloquentia, ēloquentiae} — \textquotedblleft{}eloquence\textquotedblright{}.

\begin{quote}
Ēloquentia eius iūdicēs mōvit.
\end{quote}

\begin{quote}
\emph{His eloquence moved the judges.}
\end{quote}

\emph{ēloquentia, ēloquentiae} f. Built on \emph{ēloquēns}, speaking out well. English \emph{eloquence}.

\begin{cousinweb}
\emph{ēloquentia} is built on \emph{ēloquēns} \textquotedblleft{}speaking out well\textquotedblright{} (EtymDB; Lewis \& Short), the participle of a verb made from \emph{ē} \textquotedblleft{}out\textquotedblright{} and \emph{loquor} \textquotedblleft{}to speak\textquotedblright{}. English \emph{eloquence} comes from it.
\end{cousinweb}

\begin{grammarlens}[title={What you've built}]
One more word for telling, explaining and describing at B1.
\end{grammarlens}

\subsection*{Your turn}
\begin{itemize}
\raggedright
  \item \emph{Say it:} \emph{ēloquentia, ēloquentiae}
  \item \emph{Say it:} \emph{ēloquentia, ēloquentiae}, once more
  \item \emph{Say it:} a sentence of your own with \emph{ēloquentia, ēloquentiae}
  \item \emph{Recall:} say the Latin for as, then the Latin for a speaker, then say \emph{ēloquentia, ēloquentiae} again
\end{itemize}

\begin{tcolorbox}[breakable,colback=teal!4,colframe=teal!35!black,title={Before you move on}]
What does ēloquentia, ēloquentiae mean? (\textquotedblleft{}Eloquence\textquotedblright{}.) Say it once more.
\end{tcolorbox}
\section[disputātiō, disputātiōnis]{disputātiō, disputātiōnis — a discussion, a debate}

\label{lesson:LA-C368-disputatio}



\begin{quote}
Before the new one: say the Latin for a speaker, then the Latin for eloquence.
\end{quote}

\subsection*{You'll want to know: disputātiō, disputātiōnis}
\textbf{disputātiō, disputātiōnis} — \textquotedblleft{}a discussion, a debate\textquotedblright{}.

\begin{quote}
Disputātiō dē lēgibus longa fuit.
\end{quote}

\begin{quote}
\emph{The debate about the laws was long.}
\end{quote}

\emph{disputātiō, disputātiōnis} f., the noun of \emph{disputō}, to discuss. English \emph{disputation}.

\begin{cousinweb}
\emph{disputātiō} is the noun built on \emph{disputō} \textquotedblleft{}to discuss, to argue\textquotedblright{} (EtymDB; Lewis \& Short). English \emph{disputation} comes from it.
\end{cousinweb}

\begin{grammarlens}[title={What you've built}]
One more word for telling, explaining and describing at B1.
\end{grammarlens}

\subsection*{Your turn}
\begin{itemize}
\raggedright
  \item \emph{Say it:} \emph{disputātiō, disputātiōnis}
  \item \emph{Say it:} \emph{disputātiō, disputātiōnis}, once more
  \item \emph{Say it:} a sentence of your own with \emph{disputātiō, disputātiōnis}
  \item \emph{Recall:} say the Latin for a speaker, then the Latin for eloquence, then say \emph{disputātiō, disputātiōnis} again
\end{itemize}

\begin{tcolorbox}[breakable,colback=teal!4,colframe=teal!35!black,title={Before you move on}]
What does disputātiō, disputātiōnis mean? (\textquotedblleft{}A discussion, a debate\textquotedblright{}.) Say it once more.
\end{tcolorbox}
\section[disserō, disserere]{disserō, disserere — to discuss, to hold forth (on)}

\label{lesson:LA-C368-dissero}



\begin{quote}
Before the new one: say the Latin for eloquence, then the Latin for a discussion.
\end{quote}

\subsection*{You'll want to know: disserō, disserere}
\textbf{disserō, disserere} — \textquotedblleft{}to discuss, to hold forth (on)\textquotedblright{}.

\begin{quote}
Senex dē amīcitiā disseruit.
\end{quote}

\begin{quote}
\emph{The old man spoke at length about friendship.}
\end{quote}

\emph{disserō, disserere, disseruī, dissertum}. \emph{Dē} + ablative names the subject. It is a different verb from \emph{serō}, to sow. English \emph{dissertation}.

\begin{cousinweb}
\emph{disserō} is built from \emph{dis-} \textquotedblleft{}apart\textquotedblright{} and \emph{serō, serere, seruī} \textquotedblleft{}to join, to string together\textquotedblright{} (Lewis \& Short), a different verb from \emph{serō} \textquotedblleft{}to sow\textquotedblright{}: to set ideas out in a row. English \emph{dissertation} comes from it.
\end{cousinweb}

\begin{grammarlens}[title={What you've built}]
One more word for telling, explaining and describing at B1.
\end{grammarlens}

\subsection*{Your turn}
\begin{itemize}
\raggedright
  \item \emph{Say it:} \emph{disserō, disserere}
  \item \emph{Say it:} \emph{disserō, disserere}, once more
  \item \emph{Say it:} a sentence of your own with \emph{disserō, disserere}
  \item \emph{Recall:} say the Latin for eloquence, then the Latin for a discussion, then say \emph{disserō, disserere} again
\end{itemize}

\begin{tcolorbox}[breakable,colback=teal!4,colframe=teal!35!black,title={Before you move on}]
What does disserō, disserere mean? (\textquotedblleft{}To discuss, to hold forth (on)\textquotedblright{}.) Say it once more.
\end{tcolorbox}
\section[ēloquor, ēloquī]{ēloquor, ēloquī — to speak out, to utter}

\label{lesson:LA-C368-eloquor}



\begin{quote}
Before the new one: say the Latin for a discussion, then the Latin for to discuss.
\end{quote}

\subsection*{You'll want to know: ēloquor, ēloquī}
\textbf{ēloquor, ēloquī} — \textquotedblleft{}to speak out, to utter\textquotedblright{}.

\begin{quote}
Tandem quod sentiēbat ēlocūtus est.
\end{quote}

\begin{quote}
\emph{At last he spoke out what he felt.}
\end{quote}

Deponent: \emph{ēloquor, ēloquī, ēlocūtus sum}. \emph{Ē} + \emph{loquor}: to speak out. English \emph{eloquent}.

\begin{cousinweb}
\emph{ēloquor} is built from \emph{ē} \textquotedblleft{}out\textquotedblright{} and \emph{loquor} \textquotedblleft{}to speak\textquotedblright{} (Lewis \& Short). English \emph{eloquent} comes from its present participle.
\end{cousinweb}

\begin{grammarlens}[title={What you've built}]
One more word for telling, explaining and describing at B1.
\end{grammarlens}

\subsection*{Your turn}
\begin{itemize}
\raggedright
  \item \emph{Say it:} \emph{ēloquor, ēloquī}
  \item \emph{Say it:} \emph{ēloquor, ēloquī}, once more
  \item \emph{Say it:} a sentence of your own with \emph{ēloquor, ēloquī}
  \item \emph{Recall:} say the Latin for a discussion, then the Latin for to discuss, then say \emph{ēloquor, ēloquī} again
\end{itemize}

\begin{tcolorbox}[breakable,colback=teal!4,colframe=teal!35!black,title={Before you move on}]
What does ēloquor, ēloquī mean? (\textquotedblleft{}To speak out, to utter\textquotedblright{}.) Say it once more.
\end{tcolorbox}
